\chapter{A reusable diagnostic playbook}
\label{ch:diagnostic-playbook}
\minitoc\label{a-reusable-diagnostic-playbook}
\index{diagnostic playbook}\index{find-file-literally}\index{hook bisection}\index{performance!diagnosis}

% ── Opening signpost ──────────────────────────────────────────────
This chapter distils the profiling techniques from
Chapter~\ref{ch:profiling} into a step-by-step decision tree
that draws on the GC diagnostics of Chapter~\ref{ch:gc-tuning} and
the runtime knobs of Chapter~\ref{ch:general-performance}.
Follow the steps in order whenever something is slow: each step
isolates a different layer, narrowing the search space until you
reach the root cause.

Follow this checklist in order when something is slow. Each step isolates a different layer.

\section{Step 1: Isolate. Is it Emacs or the file?}\label{step-1-isolate-is-it-emacs-or-the-file}

\begin{minted}{elisp}
;; Open the file without any mode processing:
(find-file-literally "/path/to/slow-file")
\end{minted}

If this is fast, the problem is in modes/hooks (proceed to step 3). If this is slow, the problem is I/O or GC (proceed to step 2).

\section{Step 2: Check GC health}\label{step-2-check-gc-health}

\begin{minted}{elisp}
(message "threshold: %s, gcs: %d, elapsed: %.3fs"
         gc-cons-threshold gcs-done gc-elapsed)
\end{minted}

If \texttt{gc-cons-threshold} is 800,000 or \texttt{gcs-done} is in the hundreds, something has reset your threshold. Find it:

\begin{minted}{elisp}
;; Watch for the culprit:
(add-variable-watcher
 'gc-cons-threshold
 (lambda (sym val op where)
   (when (< val (* 100 1024 1024))  ;; alert if set below 100 MB
     (message "gc-cons-threshold set to %s by: %s"
              val (or where "unknown")))))
\end{minted}

\section{Step 3: Hook bisection}\label{step-3-hook-bisection}

Find which hook function is slow:

\begin{minted}{elisp}
;; List all functions on the relevant hook:
(mapcar (lambda (fn)
          (let ((start (current-time)))
            (funcall fn)
            (cons fn (* 1000 (float-time
                              (time-subtract (current-time) start))))))
        find-file-hook)
\end{minted}

For startup hooks, dump them at the right moment:

\begin{minted}{elisp}
(add-hook 'emacs-startup-hook
          (lambda ()
            (message "startup-hook fns: %S"
                     (mapcar #'symbol-name emacs-startup-hook)))
          -99)  ;; priority -99 = run first
\end{minted}

\subsection{Cumulative hook restoration (\texttt{find-file-hook} example)}\label{cumulative-hook-restoration-find-file-hook-example}

Once you have a baseline with all hooks nil (Section~\ref{isolating-file-open-cost-nil-all-hooks-add-back-cumulatively}), add them back in groups to measure their cumulative cost. This emacsclient block restores the full \texttt{find-file-hook} and measures:

\begin{minted}{bash}
timeout 15 emacsclient -e '
(progn
  (add-hook (quote find-file-hook) (quote doom-modeline-update-persp-name))
  (add-hook (quote find-file-hook) (quote doom-modeline-update-vcs))
  (add-hook (quote find-file-hook) (quote doom-modeline-update-buffer-file-name))
  (add-hook (quote find-file-hook) (quote doom-modeline-update-buffer-file-icon))
  (add-hook (quote find-file-hook) (quote doom-modeline-update-env))
  (add-hook (quote find-file-hook) (quote recentf-track-opened-file))
  (add-hook (quote find-file-hook) (quote auto-revert--global-adopt-current-buffer))
  (add-hook (quote find-file-hook) (quote auto-revert-find-file-function))
  (add-hook (quote find-file-hook) (quote vc-refresh-state))
  (add-hook (quote find-file-hook) (quote epa-file-find-file-hook))
  (add-hook (quote find-file-hook) (quote save-place-find-file-hook))
  (let ((start (float-time)))
    (find-file-noselect "/tmp/tiny-test-hooks.txt")
    (let ((elapsed (- (float-time) start)))
      (kill-buffer "tiny-test-hooks.txt")
      (cons "+find-file-hook-all" elapsed))))' 2>&1
\end{minted}

Binary-search by commenting out half the \texttt{add-hook} calls at a time. Within 4--5 rounds you will have isolated the single most expensive entry.

\section{Step 4: \texttt{require} tracing}\label{step-4-require-tracing}

Use the timing advice from Section~\ref{wrapping-require-load-with-timing-advice}. Load it as an \texttt{-l} file:

\begin{minted}{bash}
gtimeout 15 emacs -l /path/to/require-profiler.el \
  --eval "(kill-emacs)"
\end{minted}

Sort by elapsed time. Anything over 10 ms is worth investigating.

\section{Step 5: Trace the caller}\label{step-5-trace-the-caller}

Once you know \textbf{what} is slow (e.g., \texttt{dbus} takes 17 ms), find \textbf{who} loads it using the caller tracing advice from Section~\ref{tracing-who-loads-what}.

\section{Step 5b: Binary-bisect \texttt{init.el}}\label{step-5b-binary-bisect-init.el}

If \texttt{require} tracing shows no obvious single bottleneck (costs are spread across many small packages), bisect \texttt{init.el} by loading only the first half:

\begin{minted}{bash}
head -600 ~/.emacs.d/init.el > /tmp/init-first-half.el
timeout 60 /Applications/Emacs.app/Contents/MacOS/Emacs --batch \
  -l ~/.emacs.d/early-init.el -l /tmp/init-first-half.el --eval '
(progn
  (setq org-mode-hook nil text-mode-hook nil
        outline-mode-hook nil after-change-major-mode-hook nil
        find-file-hook nil)
  (message "start: %s" (float-time))
  (find-file-noselect "~/Documents/todo/todo.org")
  (message "end: %s" (float-time)))' 2>&1 | grep -E '(start:|end:)'
\end{minted}

If the first half is fast, the bottleneck is in lines 601--end. Take 900 lines and repeat. Within 3--4 iterations you will have narrowed the culprit to a \(\sim\)50-line window.

Loading a truncated \texttt{init.el} will produce errors for missing variables. Either tolerate them (the timing is still valid) or pipe stderr away:

\begin{minted}{bash}
head -600 ~/.emacs.d/init.el > /tmp/init-first-half.el
timeout 60 /Applications/Emacs.app/Contents/MacOS/Emacs --batch \
  -l ~/.emacs.d/early-init.el -l /tmp/init-first-half.el \
  --eval "(kill-emacs)" 2>/dev/null; echo "exit: $?"
\end{minted}

A non-zero exit code means a fatal load error in that slice; reduce the line count to avoid the error-producing declaration.

\section{Step 6: Measure → Fix → Verify}\label{step-6-measure-fix-verify}

\begin{enumerate}
\item Run 5 benchmark iterations. Record the median.
\item Apply the fix.
\item Run 5 more iterations. Compare medians.
\item If the improvement is within noise (±3 ms), the fix is not performance-relevant (but may still be a correctness fix).
\end{enumerate}

\section{Step 7: Document}\label{step-7-document}

Append the fix to your fix log with:

\begin{itemize}
\item Root cause
\item Change applied (before/after code)
\item Why it works (mechanistic)
\item Regression test
\end{itemize}

The documentation is insurance against re-introducing the problem during future refactoring.

% ── Mental model summary ──────────────────────────────────────────
\section{Mental model}
Every investigation should begin with \texttt{find-file-literally},
which bypasses all mode and hook processing and therefore cleanly
separates raw file I/O cost from everything Emacs layers on top.  The
next check is GC health: if \texttt{gc-cons-threshold} has been
silently reset by a package, every measurement taken afterward is
polluted.  Once those baselines are sound, hook bisection (nilling an
entire hook list, then re-adding functions in halves) converges on the
expensive function in four or five rounds.  For load-time costs,
\texttt{require} tracing with timing advice identifies the expensive
libraries (sort by elapsed time), and capturing
\texttt{load-file-name} inside that advice answers \emph{who} triggered
each slow \texttt{require}.  When the overhead is spread thinly across
many small packages, binary-bisecting \texttt{init.el} narrows the
window efficiently.  Whatever technique you apply, always measure at
least five times, take the median, and document the root cause so the
fix survives future refactoring.
